\subsection{Sutvarkyti duomenys}
\label{subsec:sutvarkyti_duomenys}

Pradiniai duomenys buvo sutvarkyti tokiu būdu:
\begin{itemize}
    \item Klasės $2$ ir $4$ buvo atitinkamai pakeistos į $0$ ir $1$.
    \item Ištrintas ID stulpelis.
    \item Pašalintos eilutės, kuriose buvo bent viena praleista reikšmė.
    \item Duomenų eilutės buvo permaišytos atsitiktine tvarka.
    \item Požymiai buvo normalizuoti į intervalą $[0; 1]$ naudojant min-max metodą~\eqref{eq:min-max}.
\end{itemize}

Po sutvarkymo liko $683$ įrašai su $9$ požymiais ir $1$ klasės stulpeliu. Iš jų $444$ priklauso klasei $0$ (nepiktybinis navikas) ir $239$ – klasei $1$ (piktybinis navikas). Duomenų pavyzdį galima pamatyti \ref{tab:sutvarkyti_duomenys} lentelėje.

\begin{table}[H]
    \centering
    \caption{Sutvarkytų duomenų pavyzdys (pirmieji 10 įrašų).}
    \begin{tabular}{|c|c|c|c|c|c|c|c|c|c|}
        \hline
        $x_1$ & $x_2$ & $x_3$ & $x_4$ & $x_5$ & $x_6$ & $x_7$ & $x_8$ & $x_9$ & Klasė \\
        \hline
        $0{,}778$ & $0{,}222$ & $0{,}333$ & $0{,}889$ & $0{,}222$ & $1{,}000$ & $0{,}222$ & $0{,}222$ & $0{,}000$ & $1$ \\ \hline
        $0{,}778$ & $0{,}778$ & $0{,}667$ & $0{,}333$ & $1{,}000$ & $1{,}000$ & $0{,}667$ & $0{,}778$ & $0{,}667$ & $1$ \\ \hline
        $0{,}000$ & $0{,}000$ & $0{,}000$ & $0{,}000$ & $0{,}111$ & $0{,}000$ & $0{,}222$ & $0{,}000$ & $0{,}000$ & $0$ \\ \hline
        $0{,}000$ & $0{,}000$ & $0{,}111$ & $0{,}111$ & $0{,}111$ & $0{,}000$ & $0{,}222$ & $0{,}000$ & $0{,}000$ & $0$ \\ \hline
        $0{,}000$ & $0{,}000$ & $0{,}000$ & $0{,}000$ & $0{,}000$ & $0{,}000$ & $0{,}111$ & $0{,}000$ & $0{,}000$ & $0$ \\ \hline
        $0{,}444$ & $0{,}000$ & $0{,}000$ & $0{,}000$ & $0{,}222$ & $0{,}111$ & $0{,}111$ & $0{,}111$ & $0{,}000$ & $0$ \\ \hline
        $0{,}000$ & $0{,}000$ & $0{,}222$ & $0{,}000$ & $0{,}000$ & $0{,}000$ & $0{,}111$ & $0{,}000$ & $0{,}000$ & $0$ \\ \hline
        $0{,}000$ & $0{,}000$ & $0{,}000$ & $0{,}111$ & $0{,}000$ & $0{,}222$ & $0{,}000$ & $0{,}000$ & $0{,}667$ & $0$ \\ \hline
        $0{,}333$ & $0{,}111$ & $0{,}222$ & $0{,}444$ & $0{,}222$ & $0{,}778$ & $0{,}667$ & $0{,}556$ & $0{,}000$ & $1$ \\ \hline
        $0{,}333$ & $0{,}000$ & $0{,}111$ & $0{,}000$ & $0{,}111$ & $0{,}000$ & $0{,}111$ & $0{,}000$ & $0{,}000$ & $0$ \\ \hline
    \end{tabular}
    \label{tab:sutvarkyti_duomenys}
\end{table}